\documentclass[12pt]{article}


\usepackage{amsmath,amssymb,amsfonts}
\usepackage{amsthm}
\usepackage{bm}
\usepackage{geometry}
\usepackage[numbers,super,sort&compress]{natbib}
\usepackage{booktabs}
\usepackage{multirow}
\geometry{margin=1in}


% Hyperlinks (invisible style: normal text color; no boxes)
\usepackage[hidelinks]{hyperref}
\hypersetup{linktoc=all}
\usepackage[capitalize,noabbrev,nameinlink]{cleveref}


\newtheorem{definition}{Definition}[section]
\newtheorem{assumption}[definition]{Assumption}
\newtheorem{remark}[definition]{Remark}
\newtheorem{proposition}[definition]{Proposition}
\newtheorem{example}[definition]{Example}




\begin{document}


% ------------------------------------------------------------------
% STYLE GUIDELINES (first-principles, for humans and machine readers)
% - State the object and the claim once; avoid meta-commentary and repetition.
% - Lead with the operational target (what to compute / test), then give formalism.
% - Keep humility informational: scope conditions and failure modes, not performative hedging.
% - Use compact orientation aids (one toy example; short roadmaps), then move on.
% ------------------------------------------------------------------


\section{Introduction}
\label{sec:introduction}


Public discussion of automation and AI is surrounded by anxiety about job loss and wage stagnation. Important empirical questions to address include how to measure which parts of the labor market are exposed, how exposure propagates beyond an initial cluster of tasks, and which measurement choices survive falsifiable checks.


Empirical work on technological change and labor markets often relies on discrete task taxonomies---routine versus non-routine, manual versus cognitive, or coarse ``exposed'' versus ``unexposed'' categories. These partitions are useful summaries, but they do not directly encode graded similarity across tasks or spillovers from a shock that originates in one region of the task space and affects nearby tasks.


This paper develops a geometric measurement framework intended to make that structure explicit and tests whether continuous structure adds predictive value beyond discrete activity classification. We model the task structure as a metric-measure space and represent occupations as probability measures over tasks. A technology shock is represented as a task-level input profile, and its effects can propagate through a distance-based kernel. The output is a continuous occupation-level exposure index with potential neighborhood spillovers: an occupation is exposed to the extent that its task distribution lies near the shock in the task geometry.


The empirical goal is to evaluate exposure measurement in reduced form. The framework earns its keep only if the resulting exposure indices are non-redundant relative to standard discrete baselines and if the task geometry passes non-tautological diagnostics. The empirical program therefore emphasizes (i) external validation of the geometry, (ii) covariance restrictions implied by proximity and overlap, and (iii) incremental predictive content for wage and employment changes.


A central empirical finding of this paper is that continuous semantic structure significantly improves prediction of wage comovement when the kernel is properly calibrated. Normalized kernel-weighted overlap produces $t = 7.14$ (clustered standard errors), $R^2 = 0.00485$---a 191\% improvement in explanatory power over binary activity overlap ($t = 8.00$, $R^2 = 0.00167$). The semantic signal is robust to controls for occupational concentration, with overlap remaining significant ($t = 5.29$) after partialing out entropy and support size. An earlier finding that discrete structure dominates continuous was traced to an implementation artifact: bandwidth miscalibration caused kernel collapse, making all continuous measures equivalent to noise.


The empirical program has two evaluations. First, a retrospective exercise (1990--2007) compares activity-based exposure to standard industry- or occupation-mix exposure designs in settings associated with late-20th-century automation and routine-biased technological change. Second, a prospective exercise (2022--present) applies the framework to generative AI but treats the shock profile as an empirical object: multiple candidate profiles over activities are constructed and compared against rubric- or label-based baselines. This design is intended to be informative even if no single candidate profile dominates.


The paper's contribution is threefold: (i) a principled construction of occupation-activity measures from O*NET data with explicit diagnostics; (ii) a testable formalization that distinguishes continuous from discrete representations of the task domain, with methodological guidance on kernel calibration; and (iii) an empirical finding that both continuous and discrete structures are informative---continuous semantic similarity explains more variance ($R^2$) while discrete activity overlap achieves higher statistical precision ($t$-stat)---with neither representation unambiguously dominating.


Finally, wage outcomes are conditional on institutions. Productivity effects need not pass through one-for-one into wages; wage-setting and pass-through environments can introduce occupation-specific wedges. We treat this as a scope condition (\cref{subsec:wedge}) and record a ``new tasks'' extension (\cref{subsec:newtasks}), while keeping the baseline empirical focus on evaluating exposure measurement.


\cref{sec:literature} reviews related literatures. \cref{sec:theory} develops the theoretical framework, presenting both continuous and discrete representations as alternatives to be tested empirically. \cref{sec:empirics} operationalizes each component using O*NET data, reports Phase~I validation results, and specifies the Phase~II evaluation plan. \cref{sec:extensions} records extensions and historical regime signatures. \cref{sec:conclusion} concludes.


\section{Literature Review}
\label{sec:literature}


Research on automation and labor markets has increasingly adopted a task-based framework, analyzing how technologies complement or displace specific activities within occupations. In the routine-biased technological change (RBTC) paradigm, computerization disproportionately substitutes for routine, rule-based tasks while complementing non-routine analytic and interactive tasks \citep{AutorLevyMurnane2003}. Subsequent work formalized task assignment and emphasized that technological change acts on tasks, with equilibrium implications for employment reallocation and wages that cannot be recovered from a two-skill canonical model alone \citep{AcemogluAutor2011}. Empirically, this perspective underlies prominent exposure designs in which local labor markets more specialized in routine task work experienced larger subsequent declines in routine employment and a reallocation toward service work \citep{AutorDorn2013}, with cross-country evidence documenting related polarization patterns \citep{GoosManning2007}.


The limitation, for measurement, is that many task-based implementations operationalize exposure through discrete partitions (routine/non-routine; exposed/unexposed). Those partitions capture first-order patterns, but they obscure two objects that matter for reduced-form exposure regressions. First, within-category heterogeneity can be large: occupations classified as ``routine'' differ in how close their actual task bundles are to the frontier of a given technology, and ``non-routine'' occupations can still contain pockets of tasks that are technologically adjacent to the shock. Second, substitution and complementarity are naturally graded: effects propagate most strongly to nearby tasks, not only to tasks that share a label. Existing task-based exposure measures typically aggregate task-level indicators to occupations without modeling explicit spillovers across related tasks. A potential contribution of this framework is to define a metric on the task/activity domain itself and represent occupations as distributions over that domain, so that proximity and overlap in task space determine exposure propagation. Phase~I validation confirms that this continuous structure predicts outcomes better than discrete classification when properly implemented.


This concern is not purely conceptual. Work using richer task detail documents economically meaningful variation within occupations and within broad task categories \citep{AutorHandel2013,SpitzOener2006}. Related evidence in the automation-risk literature shows that occupation-level exposure measures can substantially overstate the share of employment at high risk when they ignore within-occupation task heterogeneity \citep{ArntzGregoryZierahn2016}.


Historical episodes motivate the same object, while clarifying the scope of the analogy. Electrification illustrates how a general-purpose technology can shift the demand for task bundles in heterogeneous ways, with task-based evidence in early twentieth-century manufacturing consistent with reweighting across skill-relevant activities \citep{Gray2013}. Classic accounts emphasize the emergence of technology--skill complementarity under electrification through reorganized production, capital deepening, and the rising importance of skilled maintenance and management \citep{GoldinKatz1998,GoldinKatz2008}. Technological change has heterogeneous intensity across related activities, indicating representing exposure as a continuous function of task similarity may have advantages over binary classification.


A separate literature studies exposure measurement in shift-share and related designs. The baseline empirical object is composition-predicted exposure: pre-period shares (industries, occupations, or tasks) interacted with common shocks. Work on Bartik instruments clarifies when such designs identify causal effects and how their weighting structure maps local exposure to national shocks \citep{GoldsmithPinkhamSorkinSwift2020}. More recent econometric treatments emphasize conditions under which identification can be obtained from shock variation rather than from the exogeneity of shares, along with practical guidance for inference and diagnostic checks at the shock level \citep{BorusyakHullJaravel2022}. This paper does not propose a new identification strategy; it proposes a different exposure regressor (potentially distance-propagated, task-geometry-based) and tests whether it adds content beyond standard discrete and composition-based baselines under familiar identifying assumptions.


Recent work on generative AI provides transparent, task-level occupation exposure indices that score tasks by model capability and aggregate to the occupation level \citep{EloundouEtAl2024}. Their approach operates at high task granularity, but the exposure signal is categorical (e.g., ordered classes) and aggregated via a fixed linear formula. The gap we target is different: we model exposure as a continuous function of proximity in a capability/task space and allow ``nearby'' tasks to receive partial exposure weight under an explicit kernel, rather than assigning zero weight to tasks just outside a discrete capability class. Phase~I confirms this continuous propagation mechanism improves prediction.


A growing literature on the mismatch between linguistic similarity and economic substitutability is relevant to interpreting our findings. Text embeddings capture semantic relatedness---words that appear in similar contexts---but economic substitutability depends on whether workers can actually perform different tasks, whether employers consider credentials interchangeable, and whether institutional factors constrain mobility. Despite this conceptual distinction, our empirical results show that text-embedding-based semantic similarity does predict wage comovement significantly better than discrete overlap, suggesting that linguistic similarity is at least partially informative about economic structure. The correlation between binary Jaccard overlap and semantic similarity ($r = 0.38$) indicates the measures capture related but distinct information.


Finally, measured technological exposure should be separated from distribution and pass-through. Classic and modern evidence documents that productivity growth and compensation growth can diverge for measurement and institutional reasons, including differences in deflators, changes in non-wage compensation, and shifts in labor's share \citep{BosworthPerry1994,FleckGlaserSprague2011}. More recent work argues that institutional changes affecting bargaining and rents play a first-order role in wage dynamics and in the productivity--pay relationship \citep{StansburySummers2020}. Consistent with this, evidence on automation and productivity growth shows that displacement effects can be offset in employment by downstream and general-equilibrium channels even when the labor share falls, highlighting that propagation and pass-through may differ across outcomes \citep{AutorSalomons2018}. In our framework, exposure coefficients are interpreted as reduced-form effects conditional on pass-through.


Methodologically, the framework sits alongside propagation approaches that model how shocks diffuse across related units using distance-weighted or network-based dependence structures \citep{Conley1999}. A key distinction is that many network approaches model propagation along observed linkages (e.g., input--output or supply-chain connections), whereas the kernel approach here models propagation along a latent similarity dimension in task space that must be estimated.


The construction of the task geometry is also related to text-based matching approaches that map technologies to occupations using semantic similarity between technology descriptions and job/task text \citep{Webb2020}. Longer-run adjustment through task creation remains central but is not the baseline object of estimation: automation and reinstatement via new tasks \citep{AcemogluRestrepo2018} are recorded as an extension path when the task domain must evolve.




\section{Theory}
\label{sec:theory}


This section defines the theoretical objects in continuous form: (i) the task domain as a metric-measure space; (ii) occupations as probability measures over tasks; (iii) technological shocks as task-level profiles; and (iv) propagation mechanisms that may or may not spread shocks across nearby tasks. Throughout, we develop a one-dimensional toy model to build intuition before the empirical implementation in \cref{sec:empirics}.


\subsection{The Task Manifold}
\label{subsec:task-manifold}


\begin{definition}[Task domain]
\label{def:task-domain}
The \emph{task domain} is a compact metric-measure space
\[
\mathcal{T} = (T,d,\mu),
\]
where $T$ is a compact set of technologically feasible tasks, $d : T \times T \to \mathbb{R}_{\ge 0}$ is a metric encoding reallocation frictions between tasks, and $\mu$ is a finite Borel measure on $T$ serving as a reference measure for economic mass.
\end{definition}


Distance is the primitive notion of similarity: tasks $x,y$ are close if capabilities and complements transfer at low friction. The framework uses this geometry to encode a prior on spillovers: shocks affect nearby tasks more strongly.


\begin{remark}[Manifold vs Graph]
\label{rem:graph-vs-manifold}
The manifold representation assumes continuous similarity: tasks $x, y$ at small distance $d(x,y)$ are economically interchangeable at low friction. An alternative representation treats activities as discrete nodes in a graph, where adjacency (shared occupation usage) rather than continuous distance determines relatedness. Under the graph representation, two activities are ``related'' if occupations that use one tend to use the other; continuous interpolation between activities is not meaningful. The empirical program tests which representation better predicts economic outcomes. Phase~I validation finds that both representations are informative: continuous overlap explains more variance while discrete overlap achieves higher statistical precision.
\end{remark}


\begin{example}[The Task Line]
\label{ex:task-line}
Consider a simple task domain $T = [0, 1]$ representing a single dimension of complexity (e.g., ``manual'' to ``cognitive''). The metric $d(x, y) = |x - y|$ captures the cost of retooling from task $x$ to task $y$. The reference measure $\mu$ is uniform, representing an equal density of economic value across the complexity spectrum. This one-dimensional model captures the essential structure: tasks have positions, distances matter, and mass is distributed across the space.
\end{example}


\begin{assumption}[Admissible task representation]
\label{assump:admissible-representation}
An empirical representation of the task domain is admissible if, for at least one validation target not used in construction (e.g., wage comovement, worker mobility), occupation-pair relatedness in activity space predicts stronger observed relatedness. The representation may be continuous (distance-based) or discrete (overlap-based). Specifically:
\begin{enumerate}
    \item \textbf{Metric axioms (if continuous).} The distance $\widehat d$ satisfies the metric axioms on its domain.
    \item \textbf{Scale robustness.} The representation is invariant to irrelevant rescalings of inputs (e.g., via standardization).
    \item \textbf{External diagnostic.} For at least one validation target not used in construction, greater activity-space relatedness---measured via continuous distance or discrete overlap---predicts stronger observed economic relatedness.
\end{enumerate}
\end{assumption}


The external diagnostic is a rejection criterion: it is used to rule out poor representations, not to bake success into the definition. Phase~I validation implements this diagnostic using wage comovement as the target.


\subsection{Occupations as Measures}
\label{subsec:occupations}


\begin{definition}[Occupations as measures]
\label{def:occupations}
Let $\mathcal{P}(T)$ denote the set of Borel probability measures on $T$ that are absolutely continuous with respect to $\mu$.
An occupation $j \in \{1,\dots,J\}$ is represented by a measure $\rho_j \in \mathcal{P}(T)$ with density
\[
f_j(x) = \frac{d \rho_j}{d\mu}(x),
\]
interpreted as the intensity with which occupation $j$ supplies labor to neighborhoods of task $x$.
\end{definition}


Empirically, we observe occupation-level wages and employment, not task-level outcomes. The measures $\{\rho_j\}$ define the aggregation map from task-level objects to observed outcomes and motivate diagnostics based on occupation-pair co-movement implied by proximity and overlap.


\begin{example}[Occupations on the Line]
\label{ex:occupations-line}
We model two occupations on the task line $T = [0,1]$. ``Occupation A'' (Specialist) is concentrated at the manual end, modeled as a Gaussian measure centered at $x=0.2$ with narrow variance $\sigma_A^2 = 0.01$. ``Occupation B'' (Generalist) is centered at $x=0.8$ but has wider variance $\sigma_B^2 = 0.04$, so its support extends partially into the middle range. While distinct, they share non-zero support on intermediate tasks. The overlap integral $\int \rho_A(x) \rho_B(x)\, d\mu(x)$ quantifies their task-space relatedness.
\end{example}


\subsection{Technological Fields and Shocks}
\label{subsec:tech-fields}


Technology shifts labor demand through displacement (substitution) and augmentation (complementarity). We represent both as fields over $T$.


\begin{definition}[Technological state]
\label{def:tech-state}
At time $t$, the technological state is a pair of bounded measurable functions
\[
A_t, C_t : T \to \mathbb{R},
\]
where $A_t(x)\ge 0$ is displacement intensity at task $x$ and $C_t(x)>0$ is a complementarity factor at task $x$.
\end{definition}


\begin{definition}[Technological shock profile]
\label{def:shock-field}
Let $I_t \in L^2(T,\mu)$ denote an exogenous innovation input at time $t$.
The profile $I_t(x)$ represents the intensity with which the frontier is pushed at task $x$ by R\&D, investment, or related innovation activity.
\end{definition}


In applications, $I_t$ is constructed and then potentially propagated through a measured geometry. In the prospective evaluation, multiple candidate $I_t$ are compared; the shock mapping is treated as a testable component of measurement rather than as a structural primitive. Inferring $I_t$ from labor-market outcomes is a different identification problem and is not pursued here.


We carry the complementarity field $C_t$ for notational symmetry. Baseline empirics focus on the displacement channel $(A, E^A)$; estimation of the augmentation channel $(C, E^C)$ is recorded as an extension (\cref{sec:extensions}).


\begin{example}[The Directed Shock]
\label{ex:directed-shock}
A new technology enters as a shock profile $I_t(x)$ concentrated at $x=0.25$ on the task line (e.g., a robot arm automating manual-cognitive boundary tasks). Modeled as $I_t(x) \propto \exp(-(x-0.25)^2/2\sigma_I^2)$ with $\sigma_I = 0.05$, it directly hits the core of Occupation A. However, its effect is not binary: the smooth profile assigns positive but decaying intensity to nearby tasks, rather than a hard boundary at some cutoff.
\end{example}


\subsection{Propagation of Shocks Across Tasks}
\label{subsec:propagation}


Shocks may affect only directly-targeted activities, or they may propagate to nearby activities. We present two mechanisms as alternatives to be tested empirically.


\begin{definition}[Propagation mechanisms]
\label{def:propagation-mechanisms}
\textbf{(A) Kernel propagation (continuous).} Given metric $d$ and kernel parameters $\theta$, define the normalized spillover operator $\widetilde{\mathcal{K}}_d : L^2(T,\mu) \to L^2(T,\mu)$ by
\[
[\widetilde{\mathcal{K}}_d f](x)
=
\frac{\int_T k\bigl(d(x,y); \theta\bigr)\, f(y)\, d\mu(y)}{\int_T k\bigl(d(x,y); \theta\bigr)\, d\mu(y)},
\]
where $k(\cdot;\theta)$ is nonnegative and decreasing in distance.

\textbf{(B) Identity propagation (discrete).} Define identity propagation by $[\mathcal{I} f](x) = f(x)$. Under this mechanism, shocks affect only directly-targeted activities; propagation to occupations occurs through activity sharing, not through geometric proximity.
\end{definition}


Baseline implementation of kernel propagation uses an exponential kernel with a single spillover scale:
\begin{equation}
\label{eq:exp-kernel}
k\bigl(d(x,y); \sigma\bigr)
=
\exp\!\left(- \frac{d(x,y)}{\sigma}\right),
\qquad
\theta = \sigma > 0.
\end{equation}
The normalization ensures the operator is a local averaging operator whose output scale is determined by $f$, not by $\mu$.


\begin{remark}[Selection criterion]
\label{rem:selection-criterion}
The choice between continuous and discrete propagation is empirical. If kernel-weighted overlap predicts outcomes better than binary overlap, the continuous mechanism adds value. If not, the simpler discrete mechanism is preferred. Phase~I validation (\cref{subsec:phase1}) tests this by comparing kernel-weighted occupation-pair overlap against binary Jaccard overlap as predictors of wage comovement. The result: both mechanisms are informative---continuous overlap explains more variance ($R^2$) while discrete overlap achieves higher statistical precision ($t$-stat).
\end{remark}


\begin{remark}[Bandwidth calibration]
\label{rem:bandwidth-calibration}
Kernel family (exponential vs Gaussian) and bandwidth $\sigma$ are tuning dimensions. A critical methodological finding is that $\sigma$ must be calibrated to the \emph{local} distance structure (median nearest-neighbor distance), not the global distance distribution (median of all pairwise distances). When distances cluster in a narrow band, calibrating $\sigma$ to the global median causes kernel collapse: all weights become approximately equal, destroying the signal. Calibrating to nearest-neighbor distances ensures the kernel discriminates between close and distant activity pairs.
\end{remark}


Kernel family (exponential vs Gaussian) and bandwidth $\sigma$ are treated as tuning dimensions evaluated via out-of-sample prediction, not derived from theory.


\begin{example}[Distance Decay]
\label{ex:distance-decay}
The shock from \cref{ex:directed-shock} propagates via the kernel $k(d) = e^{-d/\sigma}$ with $\sigma = 0.1$. Even though the shock centroid ($0.25$) is far from Occupation B's center ($0.8$), the ``tail'' of Occupation B's measure extends into the affected region. A binary ``exposed/unexposed'' classifier would mark Occupation B as ``0,'' but the integral
\[
E^A_B = \int_T [\widetilde{\mathcal{K}}_d I_t](x)\, d\rho_B(x)
\]
captures this tail risk as a positive but small exposure. This creates the continuous variation we exploit empirically: occupations are not simply ``in'' or ``out,'' but have graded exposure determined by the geometry of their task distributions relative to the shock.
\end{example}


\begin{remark}[Spillover decay]
\label{rem:spillover-decay}
Under the exponential kernel, $[\widetilde{\mathcal{K}}_d I](x)$ is decreasing in $d(x, \mathrm{supp}(I))$ for fixed $I$ and $\sigma$. This is immediate from the kernel definition and motivates the interpretation of $\sigma$ as a ``spillover range'' parameter. Phase~I confirms that this decay structure improves prediction beyond binary activity inclusion when properly calibrated.
\end{remark}


\begin{definition}[Baseline exposure construction---static]
\label{def:baseline-exposure}
The displacement field at time $t$ is
\[
A_t := \mathcal{P} I_t,
\]
where $\mathcal{P}$ is either the kernel operator $\widetilde{\mathcal{K}}_d$ (continuous propagation) or the identity operator $\mathcal{I}$ (discrete propagation), and $I_t$ is the exogenous shock profile over tasks. For time-aggregated exercises (e.g., retrospective 1990--2007), define $A := \mathcal{P} I$ with a fixed profile $I$ representing the cumulative shock direction.
\end{definition}


\subsection{Aggregation to Occupations and Outcome Equations}
\label{subsec:aggregation}


\begin{definition}[Occupation-level exposure functionals]
\label{def:exposure-functionals}
Given $(A_t,C_t)$ and measures $\{\rho_j\}_{j=1}^J$, define
\[
E^{A}_{j,t} \equiv \int_T A_t(x)\, d\rho_j(x),
\qquad
E^{C}_{j,t} \equiv \int_T \log C_t(x)\, d\rho_j(x).
\]
The pipeline is: (i) shock profile $I_t$ on tasks, (ii) propagation $A_t = \mathcal{P} I_t$ across tasks (continuous or discrete), (iii) aggregation via $\rho_j$ to occupation-level exposure.
\end{definition}


\begin{definition}[Discrete exposure]
\label{def:discrete-exposure}
Under identity propagation ($\mathcal{P} = \mathcal{I}$), occupation-level exposure reduces to:
\[
E^A_j = \sum_{a \in T_n} I(a) \cdot \rho_j(a),
\]
which is a weighted sum of shock intensity over the occupation's activity distribution. This is equivalent to:
\[
E^A_j = \sum_{a \in T_n} I(a) \cdot \mathbb{1}[\rho_j(a) > 0] \cdot \rho_j(a),
\]
where $\mathbb{1}[\cdot]$ is the indicator function. Exposure depends only on which activities the occupation uses and how intensively it uses them, not on distances between activities.
\end{definition}


\begin{definition}[Reduced-form outcome equations (partial equilibrium)]
\label{def:reduced-form}
Let $w_{j,t}$ and $L_{j,t}$ denote wages and employment. We evaluate
\[
\Delta \log w_{j,t} = \beta_w E^{A}_{j,t} + \gamma_w E^{C}_{j,t} + \bm{X}_{j,t}'\delta_w + u_{j,t},
\]
\[
\Delta \log L_{j,t} = \beta_L E^{A}_{j,t} + \gamma_L E^{C}_{j,t} + \bm{X}_{j,t}'\delta_L + v_{j,t},
\]
interpreting coefficients as conditional associations used to compare exposure measures.
\end{definition}


\begin{remark}[Partial-equilibrium scope]
\label{rem:pe-scope}
The reduced form abstracts from general-equilibrium feedbacks and re-sorting. This paper uses \cref{def:reduced-form} to evaluate whether the geometric exposure indices contain incremental information relative to baselines; it does not interpret coefficients structurally.
\end{remark}


\begin{remark}[Implied covariance structure---diagnostic target]
\label{rem:covariance}
If outcomes load on exposure with common sign and shocks are spatially concentrated in task space, then occupation pairs with greater task-distribution overlap should exhibit stronger outcome comovement.

\textbf{Kernel-weighted overlap.} Define occupation-pair proximity as
\[
\mathrm{Overlap}_K(i,j) = \int_T \int_T \rho_i(x)\, \rho_j(y)\, k(d(x,y); \sigma)\, d\mu(x)\, d\mu(y).
\]

\textbf{Normalized overlap.} To control for concentration effects (specialist vs generalist occupations), define normalized overlap:
\[
\mathrm{NormOverlap}_K(i,j) = \frac{\mathrm{Overlap}_K(i,j)}{\sqrt{\mathrm{Overlap}_K(i,i) \cdot \mathrm{Overlap}_K(j,j)}}.
\]
This measures the angular similarity between occupation distributions in kernel space, removing scale effects from concentration.

\textbf{Binary overlap.} Define binary overlap as the Jaccard index:
\[
\mathrm{Overlap}_B(i,j) = \frac{|A_i \cap A_j|}{|A_i \cup A_j|},
\]
where $A_j = \{a : \rho_j(a) > 0\}$ is the activity support of occupation $j$.

Phase~I tests whether kernel-weighted overlap improves prediction of wage comovement beyond binary overlap.
\end{remark}


\subsection{Why Continuous Structure Might Dominate}
\label{subsec:continuous}


The manifold representation assumes that economic frictions vary smoothly with task distance. Under this model, occupations with similar task distributions experience correlated shocks even if they do not share identical activities.


\begin{remark}[Semantic similarity as economic signal]
\label{rem:semantic-signal}
Text embeddings capture semantic relatedness---activities described in similar language tend to require related capabilities. If labor market frictions depend on capability transferability, semantic similarity should predict economic outcomes. Phase~I confirms this: semantic overlap predicts wage comovement significantly better than binary activity overlap, suggesting that ``nearby but not identical'' activities are economically relevant.
\end{remark}


\begin{remark}[Concentration as noise]
\label{rem:concentration-noise}
Occupation measures vary in concentration: specialists have mass concentrated on few activities, while generalists spread mass across many. Unnormalized overlap conflates semantic similarity with concentration similarity. The normalized overlap measure removes this confound. Empirically, normalization \emph{improves} prediction (R$^2$ increases by 56.6\%), indicating that concentration effects were noise obscuring the semantic signal, not signal themselves.
\end{remark}


\subsection{Why Discrete Structure Might Dominate}
\label{subsec:threshold}


An alternative microfoundation, consistent with discrete dominance, models hiring and production as threshold-based processes.


\begin{definition}[Threshold production]
\label{def:threshold-production}
Occupation $j$ produces output using activities in its support $A_j = \{a : \rho_j(a) > 0\}$. For each activity $a \in A_j$, the worker either can perform $a$ (contributing to output) or cannot (contributing zero). There is no ``partial credit'' for performing a similar activity $a'$ with small $d(a, a')$.
\end{definition}


Under this model, a technology shock to activity $a$ affects occupation $j$ if and only if $a \in A_j$. The kernel-weighted propagation mechanism, which assigns positive exposure to occupations with $a' \approx a$ in their support, would overstate exposure by treating ``nearby but not identical'' as economically relevant.


\begin{remark}[Hiring as checkbox verification]
\label{rem:checkbox}
The threshold interpretation is consistent with one view of how hiring works: employers verify whether candidates can perform specific required tasks, not whether they can perform ``similar'' tasks. A job posting for a Python developer requires Python proficiency; proficiency in a ``nearby'' language (say, Ruby) does not substitute, even if the semantic distance $d(\text{Python}, \text{Ruby})$ is small.
\end{remark}


\begin{remark}[Empirical test of threshold model]
\label{rem:threshold-test}
The threshold model predicts: (i) binary activity overlap should predict wage comovement, (ii) continuous distance weighting should not improve prediction beyond binary overlap, and (iii) the spillover bandwidth $\sigma$ should be inert. Phase~I partially falsifies this: continuous overlap explains significantly more variance than binary overlap when properly calibrated, and $\sigma$ matters. However, binary overlap achieves higher statistical precision ($t = 8.00$ vs $t = 7.14$ clustered), suggesting the threshold model captures a robust but coarser signal. Both representations are informative.
\end{remark}


\subsection{Wage-setting wedge and pass-through (scope condition)}
\label{subsec:wedge}


Observed wages need not track marginal products one-for-one. We record this as a scope condition rather than a mechanism to be estimated.


\begin{definition}[Occupation-level wage-setting wedge]
\label{def:wedge}
Let $\mathrm{MPL}_{j,t}$ denote the (unobserved) marginal product of labor for occupation $j$. Define $b_{j,t} > 0$ such that
\[
\log w_{j,t} = \log \mathrm{MPL}_{j,t} + \log b_{j,t}.
\]
\end{definition}


Empirically, proxies for wage-setting environments (e.g., union coverage, licensing prevalence, concentration) enter as controls or descriptive correlates. Interpretation of exposure coefficients is conditional on these environments \citep{StansburySummers2020,AutorSalomons2018}.


\subsection{New tasks and an evolving task domain (extension)}
\label{subsec:newtasks}


The baseline holds $(T,\mu)$ fixed. Over longer horizons, outcomes depend on task creation and reweighting \citep{AcemogluRestrepo2018}. This section records extension paths: expanding $T_t$, time-varying $\mu_t$, and time-varying occupation measures $\rho_{j,t}$.


Empirically, distinguishing new tasks from broad augmentation is difficult with occupation-level outcomes alone. Higher-frequency text embeddings are the natural route for identifying emerging task coordinates over time.




\section{Empirical Strategy}
\label{sec:empirics}


This section operationalizes the theoretical framework using O*NET data. We describe how to (i) construct a discrete approximation of the task manifold, (ii) build occupation measures from ratings, (iii) compute activity distances, and (iv) evaluate the resulting exposures against baselines. The tests are organized to separate geometry validation (Phase I) from shock application and exposure comparison (Phase II).


\subsection{Discrete Approximation of the Manifold}
\label{subsec:discrete-approx}


The theoretical objects of \cref{sec:theory} are instantiated on a finite activity domain derived from O*NET.


\begin{definition}[Empirical activity domain]
\label{def:empirical-activity-domain}
Let $T_n = \{a_1, \ldots, a_n\}$ be a finite set of work activities. In baseline implementation:
\begin{itemize}
    \item $T_n$ is the set of O*NET Detailed Work Activities (DWA; $n = 2{,}087$), with Generalized Work Activities (GWA; $n \approx 41$) as a robustness domain \citep{onet2024database,onet2024content}.
    \item The reference measure $\mu$ is a probability mass function on $T_n$. Baseline: uniform ($\mu(a) = 1/n$ for all $a$).
    \item All integrals $\int_T f(x)\, d\mu(x)$ reduce to finite sums $\sum_{a \in T_n} f(a)\, \mu(a)$.
\end{itemize}
The metric $d: T_n \times T_n \to \mathbb{R}_{\ge 0}$ is constructed empirically (\cref{subsubsec:recipes}).
\end{definition}


\begin{remark}[Constructing occupation measures from O*NET]
\label{rem:rho-construction}
The occupation measure $\rho_j$ is constructed from O*NET ratings via the following steps \citep{onet2024dwa,onet2024scales}:


\medskip
\noindent\textbf{(A) Raw weight construction $w_j(a)$:}


\emph{If $T_n = $ GWA (Work Activities file):}
\begin{itemize}
    \item Retrieve Importance $\mathrm{IM}_{j,a} \in [1,5]$ and Level $\mathrm{LV}_{j,a} \in [0,7]$ for each occupation-activity pair.
    \item Baseline: $w_j(a) = \mathrm{IM}_{j,a} \times \mathrm{LV}_{j,a}$ after rescaling both to $[0,1]$.
    \item Robustness variants: IM only; LV only; first principal component of (IM, LV).
\end{itemize}


\emph{If $T_n = $ DWA (Detailed Work Activities):}
\begin{itemize}
    \item DWAs lack direct importance ratings; importance is derived via linked tasks.
    \item Let $\mathrm{Tasks}(j, a) = \{\text{tasks in occupation } j \text{ linked to DWA } a\}$.
    \item Baseline aggregator: $w_j(a) = \max_{t \in \mathrm{Tasks}(j,a)} \mathrm{IM}_{j,t}$ (top importance).
    \item Robustness aggregators: mean importance; relevance-weighted mean; top importance among ``core tasks'' only (core $=$ top 70\% by task importance within occupation, pre-committed threshold).
    \item The ``top importance'' rule follows O*NET's published DWA ranking methodology.
\end{itemize}


\medskip
\noindent\textbf{(B) Suppression and missingness rules:}
\begin{itemize}
    \item Set $w_j(a) = 0$ if O*NET flags ``Recommend Suppress $=$ Y'' for that cell.
    \item Set $w_j(a) = 0$ if ``Not Relevant $=$ Y'' (per O*NET's 75\% rule for Level).
    \item If an occupation has fewer than 10 activities with nonzero weight after filtering, flag the occupation and report coverage in diagnostics.
    \item Imputation: none in baseline; sensitivity check with minimum-weight floor (e.g., $w = 0.01$ for suppressed cells) as robustness.
\end{itemize}


\medskip
\noindent\textbf{(C) Normalization to probability measure:}
\[
\rho_j(a) = \frac{w_j(a)}{\sum_{a' \in T_n} w_j(a')}.
\]
This yields $\rho_j \in \mathcal{P}(T_n)$. The density with respect to $\mu$ is:
\[
f_j(a) = \frac{\rho_j(a)}{\mu(a)} = n \cdot \rho_j(a) \quad \text{under uniform } \mu.
\]
Under non-uniform $\mu$ (robustness variants), use $f_j(a) = \rho_j(a)/\mu(a)$ directly.


\medskip
\noindent\textbf{(D) Implementation summary (O*NET 30.0):}

The baseline implementation uses:
\begin{itemize}
    \item Files: Task Ratings.xlsx (task importance by occupation), Tasks to DWAs.xlsx (task-to-DWA linkages), DWA Reference.xlsx (DWA definitions).
    \item Scale handling: IM scale $[1,5]$ rescaled to $[0,1]$ via $(x-1)/4$; LV scale $[0,7]$ rescaled to $[0,1]$ via $x/7$.
    \item Suppression: drop if Recommend Suppress $=$ ``Y'' or $N < 5$.
    \item Aggregation: MAX importance over linked tasks per O*NET methodology.
    \item Coverage: 894 occupations with valid measures after filtering.
\end{itemize}

\medskip
\noindent\textbf{(E) Distribution statistics:}
\begin{itemize}
    \item Sparsity: 99.1\% of occupation-activity weights are zero.
    \item Median activities per occupation: 19 (of 2,087 possible).
    \item Mean activities per occupation: 19.6 (std: 6.7).
    \item Effective support distribution: most occupations use a small fraction of the activity space.
    \item Median pairwise Jaccard overlap: 0.0 (most occupation pairs share no activities).
\end{itemize}
\end{remark}


\begin{remark}[Matrix representation of the spillover operator]
\label{rem:matrix-notation}
In discrete instantiation with uniform $\mu$, the spillover operator becomes a matrix $K \in \mathbb{R}^{n \times n}$:
\[
K_{ab} = k(d(a,b); \sigma) = \exp(-d(a,b)/\sigma).
\]
Note: we do \emph{not} row-normalize $K$ in the baseline implementation. Row normalization, while theoretically motivated for probability-preserving operators, empirically destroys the signal when $n$ is large (2,087 activities). With unnormalized $K$, the propagated field is $A = K I$ and occupation exposure is $E^A_j = \rho_j^\top A$.
\end{remark}


\begin{remark}[Discrete overlap]
\label{rem:discrete-overlap}
The occupation-pair overlap from \cref{rem:covariance} has discrete form:
\[
\mathrm{Overlap}_K(i,j) = \rho_i^\top K \rho_j = \sum_{a,b \in T_n} \rho_i(a)\, K_{ab}\, \rho_j(b).
\]
The normalized version is:
\[
\mathrm{NormOverlap}_K(i,j) = \frac{\rho_i^\top K \rho_j}{\sqrt{(\rho_i^\top K \rho_i)(\rho_j^\top K \rho_j)}}.
\]
Under identity propagation ($K = I$), these reduce to the weighted inner product $\rho_i^\top \rho_j$ and its cosine-normalized version.
\end{remark}


\subsubsection{Construction recipes for the activity metric}
\label{subsubsec:recipes}


\begin{remark}[Recipes for constructing $\widehat d$ on activities]
\label{rem:recipes}
Instantiate $T$ as a finite set of activities (GWA or DWA). Occupations are probability measures over activities, not points in $T$.
\begin{enumerate}
    \item \textbf{Recipe X (rating-cooccurrence geometry).} Represent each activity $a \in T_n$ by its importance profile across occupations: $v(a) = (w_1(a), \ldots, w_J(a))$, where $w_j(a)$ is occupation $j$'s weight on activity $a$ (defined in \cref{rem:rho-construction}). Standardize, reduce dimension (PCA retaining 90\% variance), and define $\widehat d(a,b)$ as Euclidean distance in the reduced space.
    
    \emph{Interpretation:} activities are close if they co-occur with similar importance across the occupation distribution.
    
    \emph{Validation targets:} worker mobility, wage comovement (external to O*NET ratings).
    
    \emph{Phase~I outcome:} Failed initial validation due to bandwidth miscalibration; not pursued further.
    
    \item \textbf{Recipe Y (text-embedding geometry).} Embed activity titles/descriptions (GWA definitions or DWA titles) using a sentence encoder (all-mpnet-base-v2, 768 dimensions). Define $\widehat d(a,b)$ as cosine distance in embedding space.
    
    \emph{Validation targets:} mobility, wage comovement.
    
    \emph{Phase~I outcome:} Passed validation after bandwidth correction. Normalized overlap achieves $t = 7.14$, $R^2 = 0.00485$, a 191\% improvement in $R^2$ over binary Jaccard.
    
    \item \textbf{Binary Jaccard (no distance).} Define activity-space relatedness via set intersection without any distance weighting:
    \[
    \mathrm{Jaccard}(i,j) = \frac{|A_i \cap A_j|}{|A_i \cup A_j|},
    \]
    where $A_j = \{a : \rho_j(a) > 0\}$.
    
    \emph{Phase~I outcome:} Passed validation---$t = 8.00$, $R^2 = 0.00167$. Serves as baseline for continuous measures.
\end{enumerate}
\end{remark}


\cref{tab:validation-matrix} summarizes the non-tautological principle used in Phase~I and the validation outcomes.


\begin{table}[h!]
\centering
\begin{tabular}{|l|l|l|l|}
\hline
\textbf{Distance construction} & \textbf{What it uses} & \textbf{Valid targets} & \textbf{Phase I result} \\
\hline
Recipe Y (text embedding) & activity text & mobility; wage & Passed ($t = 7.14$, $R^2 = 0.00485$) \\
\hline
Binary Jaccard & activity support & mobility; wage & Passed ($t = 8.00$, $R^2 = 0.00167$) \\
\hline
\end{tabular}
\caption{Validation targets by activity-distance recipe and Phase~I outcomes. Continuous semantic overlap explains more variance ($R^2$); binary overlap achieves higher statistical precision ($t$-stat). Both pass validation.}
\label{tab:validation-matrix}
\end{table}


\subsection{Design Choices and Pre-Commitment}
\label{subsec:design-choices}


This subsection summarizes the design axes that define the empirical pipeline. Each axis has a baseline choice and pre-committed robustness variants. All choices are fixed before examining outcome data.


\medskip
\noindent\textbf{DESIGN AXIS CHECKLIST:}


\medskip
\noindent$\square$ \textbf{Activity domain $T_n$}
\begin{itemize}
    \item Baseline: DWA (Detailed Work Activities; $n = 2{,}087$)
    \item Robustness: GWA (Generalized Work Activities; $n \approx 41$)
    \item Stress test (optional): Task statements ($n \approx 19{,}000$)
\end{itemize}


\noindent$\square$ \textbf{Reference measure $\mu$}
\begin{itemize}
    \item Baseline: uniform ($\mu(a) = 1/n$)
    \item Robustness: hierarchy-weighted (GWA mass split to child DWAs); average-importance-weighted ($\mu(a) \propto \mathrm{mean}_j\, w_j(a)$)
\end{itemize}


\noindent$\square$ \textbf{Occupation measure construction $w_j(a)$}
\begin{itemize}
    \item Baseline (GWA): IM $\times$ LV, rescaled to $[0,1]$
    \item Baseline (DWA): top importance over linked tasks
    \item Robustness: IM only; mean importance; core-task filter
\end{itemize}


\noindent$\square$ \textbf{Activity distance recipe}
\begin{itemize}
    \item Baseline: Recipe Y (text embedding with proper bandwidth calibration)
    \item Robustness: Binary Jaccard (discrete baseline)
    \item Phase~I finding: continuous semantic distance improves prediction when properly calibrated
\end{itemize}


\noindent$\square$ \textbf{Kernel specification}
\begin{itemize}
    \item Baseline: exponential, $k(d;\sigma) = \exp(-d/\sigma)$, \textbf{unnormalized}
    \item Robustness: Gaussian, $k(d;\sigma) = \exp(-d^2/2\sigma^2)$
    \item $\sigma$: median of nearest-neighbor distances (not global distance median)
    \item Phase~I finding: $\sigma$ calibrated to NN distances ($\sigma = 0.223$) yields discrimination ratio 4.5; global median ($\sigma = 0.744$) causes kernel collapse
\end{itemize}


\noindent$\square$ \textbf{Overlap normalization}
\begin{itemize}
    \item Baseline: Normalized overlap (controls for concentration)
    \item Robustness: Unnormalized overlap
    \item Phase~I finding: Normalization improves $R^2$ by 56.6\%; concentration is noise, not signal
\end{itemize}


\noindent$\square$ \textbf{Shock profiles (Phase~II)}
\begin{itemize}
    \item Primary: $v_3$ (external anchor)
    \item Robustness/hypothesis tests: $v_1$ (capability-only), $v_2$ (capability $+$ structure)
    \item Selection rule: no outcome-based selection; primary designated ex ante
\end{itemize}


\noindent$\square$ \textbf{Validation targets (Phase I)}
\begin{itemize}
    \item External monotonicity: wage comovement (OES)
    \item Baselines: binary Jaccard; routine/manual indicators; categorical AI-exposure scores
    \item Alternative target (extension): worker mobility (CPS or SIPP transitions)
\end{itemize}


\noindent$\square$ \textbf{Outcome evaluation (Phase II)}
\begin{itemize}
    \item Held-out structure: by time window (training: pre-period; test: post-period) or by region (training: subset of CZs; test: remainder)
    \item Headline metric: incremental $R^2$ vs discrete baselines
\end{itemize}


\subsection{Constructing the activity geometry and occupation measures}
\label{subsec:empirical-manifold}


Construction proceeds in four steps:


\medskip
\noindent\textbf{Step 1: Choose activity domain $T_n$.}
\begin{itemize}
    \item Baseline: DWA ($n = 2{,}087$). Rationale: DWAs are shared across occupations (enabling overlap analysis), granular enough to capture within-category heterogeneity, and have known nesting structure for hierarchy-weighted $\mu$.
    \item Robustness: GWA ($n \approx 41$) to test whether coarser domain destroys signal.
    \item Pruning: drop activities with near-zero mass across all occupations ($\sum_j w_j(a) < $ threshold; threshold pre-committed at $p5$ of activity totals).
\end{itemize}


\medskip
\noindent\textbf{Step 2: Construct occupation measures $\rho_j$.}
\begin{itemize}
    \item Apply the $w_j(a)$ construction rules (\cref{rem:rho-construction}) using baseline aggregator.
    \item Apply suppression/missingness rules.
    \item Normalize to probability measures.
    \item Output: $J \times n$ matrix of $\rho_j(a)$ values.
\end{itemize}


\medskip
\noindent\textbf{Step 3: Construct activity distance and kernel.}
\begin{itemize}
    \item Embed activity titles using all-mpnet-base-v2 (768 dimensions).
    \item Compute cosine distance matrix: $d(a,b) = 1 - \cos(\mathrm{emb}(a), \mathrm{emb}(b))$.
    \item Compute nearest-neighbor distances: for each activity $a$, find $d_{\mathrm{NN}}(a) = \min_{b \neq a} d(a,b)$.
    \item Set $\sigma = \mathrm{median}(d_{\mathrm{NN}})$. Empirically: $\sigma = 0.223$.
    \item Compute kernel matrix: $K_{ab} = \exp(-d(a,b)/\sigma)$. Do not row-normalize.
\end{itemize}


\medskip
\noindent\textbf{Step 4: Produce audit outputs.}
\begin{itemize}
    \item Nearest-neighbor table: for 20 activities spanning domains, list 5 nearest neighbors under $d$ and verify face validity.
    \item $\rho_j$ diagnostics: entropy distribution; effective support distribution; flag sparse occupations.
    \item Overlap distribution: histogram of occupation-pair overlap values.
    \item Kernel discrimination check: verify that kernel weights span a meaningful range (discrimination ratio $> 3$).
\end{itemize}


\subsection{The Diagnostics Gateway: Phase I Validation}
\label{subsec:phase1}


Phase I acts as a \emph{selection mechanism} for the activity-space representation. The continuous geometry (Recipe Y with proper calibration) was tested and passed; binary activity overlap also passed but with lower explanatory power. We report these findings as substantive empirical results.


Phase I tests whether the constructed representation is coherent and passes basic rejection diagnostics. Phase I is necessary but not sufficient: it filters out pathological constructions but does not validate that exposure indices predict outcomes.


\medskip
\noindent\textbf{DIAGNOSTIC A: $\rho_j$ coherence checks}


\medskip
\noindent\textit{(A1) Sparsity and entropy.}
\begin{itemize}
    \item Compute entropy $H(\rho_j) = -\sum_a \rho_j(a) \log \rho_j(a)$ for each occupation.
    \item Finding: The occupation-activity matrix is 99.1\% zeros. Median Jaccard overlap between occupation pairs is 0.0. Most occupations share no activities at all.
    \item Interpretation: The activity-sharing structure is extremely sparse and fragmented.
\end{itemize}


\noindent\textit{(A2) Face validity spot checks.}
\begin{itemize}
    \item Selected 10 occupations spanning sectors (nurse, electrician, accountant, truck driver, software developer, lawyer, retail cashier, machinist, teacher, manager).
    \item For each, top-5 activities by $\rho_j(a)$ were manually reviewed.
    \item Finding: Top activities match occupation descriptions; construction appears valid.
\end{itemize}


\medskip
\noindent\textbf{DIAGNOSTIC B: Kernel calibration check}


\medskip
\noindent\textit{(B1) Distance distribution.}
\begin{itemize}
    \item Activity embedding distances: mean $= 0.734$, std $= 0.132$, range $= [0.02, 1.15]$.
    \item Effective range (p95 $-$ p5) $= 0.43$.
    \item Nearest-neighbor distances: median $= 0.223$, mean $= 0.229$.
\end{itemize}


\noindent\textit{(B2) Kernel collapse diagnosis.}

Using global distance median ($\sigma = 0.744$) causes kernel collapse:
\begin{itemize}
    \item Kernel weight range (p95 $-$ p5) $= 0.00025$.
    \item All weights approximately equal ($\approx 1/2087$).
    \item Discrimination ratio $= 1.57$ (insufficient).
\end{itemize}

Using NN distance median ($\sigma = 0.223$) restores discrimination:
\begin{itemize}
    \item Discrimination ratio $= 4.54$.
    \item Kernel weights span meaningful range.
\end{itemize}


\medskip
\noindent\textbf{DIAGNOSTIC C: External monotonicity (wage comovement)}


\medskip
\noindent\textit{(C1) Regression specification.}
\[
\text{WageComovement}_{i,j} = \alpha + \beta \times \text{Overlap}_{i,j} + \varepsilon_{i,j}
\]
Standard errors clustered by origin occupation. $n = 246{,}051$ occupation pairs, $n_{\text{clusters}} = 701$.


\noindent\textit{(C2) Primary results.}

\begin{table}[h!]
\centering
\begin{tabular}{lccccc}
\toprule
\textbf{Measure} & $\beta$ & SE & $t$ & $p$ & $R^2$ \\
\midrule
Normalized kernel overlap & 0.322 & 0.045 & 7.14 & $< 10^{-12}$ & 0.00485 \\
Unnormalized kernel overlap & 1.904 & 0.323 & 5.90 & $< 10^{-8}$ & 0.00310 \\
Binary Jaccard & 0.471 & 0.059 & 8.00 & $< 10^{-15}$ & 0.00167 \\
\bottomrule
\end{tabular}
\caption{Phase I validation results: occupation-pair overlap as predictor of wage comovement. All measures use clustered standard errors. Normalized kernel overlap achieves highest $R^2$.}
\label{tab:phase1-results}
\end{table}


\noindent\textit{(C3) Robustness checks.}

\begin{table}[h!]
\centering
\begin{tabular}{lll}
\toprule
\textbf{Check} & \textbf{Result} & \textbf{Interpretation} \\
\midrule
Permutation test ($n=1000$) & $p < 0.001$ for all measures & Effects not due to chance \\
Random baseline ($n=100$) & Semantic at 100th percentile & Semantic structure is real \\
Entropy control & $t_{\text{overlap}} = 5.29$ after control & Not driven by concentration \\
Support control & $t_{\text{overlap}} = 5.43$ after control & Not driven by job breadth \\
Cross-validation (5-fold) & Overfit ratio $< 1.2$ for all & No overfitting \\
\bottomrule
\end{tabular}
\caption{Phase I robustness checks. Continuous semantic structure is robust to all tests.}
\label{tab:phase1-robustness}
\end{table}


\noindent\textit{(C4) Interpretation.}

Normalized kernel-weighted overlap explains substantially more variance in wage comovement than binary Jaccard overlap. The improvement in $R^2$ is 191\% (from 0.00167 to 0.00485). However, binary Jaccard achieves higher statistical precision ($t = 8.00$ vs $t = 7.14$ for normalized kernel), suggesting the semantic signal is noisier despite explaining more variance. This is not driven by concentration effects: after controlling for entropy and support size, the semantic overlap coefficient remains significant ($t = 5.29$).

The semantic signal is robust across validation checks:
\begin{itemize}
    \item Semantic overlap is at the 100th percentile of the random baseline distribution (100 seeds).
    \item Permutation $p < 0.001$: the effect is not due to spurious correlation.
    \item Cross-validation overfit ratio $< 1.2$: results generalize to held-out data.
\end{itemize}


\noindent\textit{(C5) Implementation artifact resolved.}

An earlier version (v0.5.0) concluded that ``discrete dominates continuous'' based on kernel-weighted overlap failing permutation testing. This was traced to bandwidth miscalibration: setting $\sigma$ to the global distance median (0.744) caused kernel collapse, making all continuous measures equivalent to noise. With proper calibration ($\sigma = 0.223$, the NN median), the continuous measure explains substantially more variance than the discrete baseline, though with lower statistical precision.


\medskip
\noindent\textbf{PHASE I SELECTION OUTCOME:}

\begin{itemize}
    \item Normalized kernel overlap passes all validation criteria with highest $R^2$.
    \item Binary Jaccard passes validation but with lower explanatory power (baseline).
    \item Continuous representation is preferred; discrete is retained as robustness check.
\end{itemize}

\textbf{Primary measure for Phase II:} Normalized kernel-weighted overlap with importance weighting:
\[
E^A_j = \frac{\sum_{a,b \in T_n} I(b) \cdot K_{ab} \cdot \rho_j(a)}{\sqrt{\sum_{a,b} K_{ab} \cdot \rho_j(a) \cdot \rho_j(b)}}.
\]


\subsection{Experiment A: Retrospective evaluation (1990--2007)}
\label{subsec:experimentA}


The retrospective evaluation compares activity-based exposure to standard composition-based exposure designs in settings associated with late-20th-century automation.


\medskip
\noindent\textbf{BASELINE EXPOSURE CONSTRUCTION:}

Given shock profile $I(a)$ over activities, occupation-level exposure is:
\[
E^A_j = \frac{\rho_j^\top K I}{\sqrt{\rho_j^\top K \rho_j}},
\]
where $K$ is the unnormalized kernel matrix and the denominator normalizes for concentration. This is a kernel-weighted, concentration-normalized sum of shock intensity over the occupation's activity distribution.


\medskip
\noindent\textbf{SHOCK PROFILE SPECIFICATION (pre-committed):}


The retrospective evaluation uses a Routine-Biased Technological Change (RBTC) shock profile $I$ over activities.


\medskip
\noindent\textit{Baseline (time-invariant):}
\begin{itemize}
    \item Define $I(a)$ based on activity characteristics: positive loading on activities classified as routine/repetitive (using O*NET ``Structured vs Unstructured Work'' or equivalent); negative loading on non-routine analytic and manual activities.
    \item Compute exposure once; interact with long-difference outcomes (1990--2007).
\end{itemize}


\noindent\textit{Robustness (time-varying):}
\begin{itemize}
    \item Allow $I_t$ to vary by decade (1990s vs 2000s) if supported by auxiliary evidence on automation timing.
    \item Cumulative exposure: $\sum_t A_t$ with equal weights across years.
\end{itemize}


\noindent Explicitly: we do not estimate $I$ from outcomes. The shock profile is constructed from activity characteristics and applied via the occupation measures.


\medskip
\noindent\textbf{ROBUSTNESS PANELS:}
\begin{itemize}
    \item Domain: GWA vs DWA
    \item $\mu$: uniform vs hierarchy-weighted
    \item Normalization: normalized vs unnormalized overlap
    \item Distance: kernel-weighted vs binary Jaccard (expected to perform worse based on Phase~I)
\end{itemize}


\medskip
\noindent\textbf{DOMINANCE TEST:}

To assess whether activity-based exposure dominates other channels, estimate:
\[
Y_j = \alpha + \beta_1 E^{\text{activity}}_j + \beta_2 E^{\text{industry}}_j + \beta_3 E^{\text{credential}}_j + \beta_4 E^{\text{wage}}_j + \epsilon_j,
\]
where $E^{\text{industry}}$ is industry-composition exposure (Bartik), $E^{\text{credential}}$ is credential-similarity exposure, and $E^{\text{wage}}$ is wage-proximity exposure.

If $\beta_1$ is not the dominant coefficient, activity-based exposure may be one of several channels rather than the primary mechanism.


\medskip
\noindent\textbf{OUTPUT:}
\begin{itemize}
    \item Comparison table: activity-based exposure vs composition-based Bartik exposure vs routine-share exposure.
    \item Dominance test: coefficients on activity-based vs alternative channels.
    \item Figure: within-group dispersion captured by activity overlap (e.g., distribution of activity-based exposure within ``routine'' category).
\end{itemize}


\subsection{Experiment B: Prospective evaluation (2022--present)}
\label{subsec:experimentB}


The prospective evaluation applies the framework to generative AI. The object is comparison among pre-specified shock profiles, not identification of a single ``true'' profile.


\medskip
\noindent\textbf{CANDIDATE SHOCK PROFILES (activity-level):}


Each candidate defines $I^{(v)}(a)$ for activities $a \in T_n$.


\medskip
\noindent\textbf{$v_1$ (capability-only):}
\begin{itemize}
    \item Construction: Score each activity on ``language/cognitive capability relevance'' using a rule-based mapping.
    \item If $T_n = $ GWA: positive loading on activities tagged as involving information processing, written comprehension, oral expression; negative loading on physical/manual activities.
    \item If $T_n = $ DWA: propagate GWA-level scores to DWAs via nesting structure ($I^{v_1}(\mathrm{dwa}) = I^{v_1}(\text{parent GWA})$).
\end{itemize}


\noindent\textbf{$v_2$ (capability $+$ structure):}
\begin{itemize}
    \item Construction: $v_1$ loadings plus additional positive weight on activities tagged as structured/repetitive.
    \item ``Structured'' operationalized as: activities under GWAs related to ``Processing Information,'' ``Documenting/Recording Information,'' or with O*NET ``Structured vs Unstructured'' rating above median.
    \item Hypothesis: structure amplifies AI exposure beyond raw capability.
\end{itemize}


\noindent\textbf{$v_3$ (external anchor---PRIMARY):}
\begin{itemize}
    \item Construction: Derive from an independent source not using O*NET geometry.
    \item Options (choose one and document):
    \begin{enumerate}
        \item[(a)] Human expert ratings: Domain experts rate activities on ``susceptibility to generative AI automation'' without seeing $d$ or $\rho_j$.
        \item[(b)] Published rubric mapping: Use Eloundou et al.\ (2024) exposure rubrics; map rubric categories to activities via text matching (rubric text $\leftrightarrow$ activity text), not via $d$.
        \item[(c)] Patent-activity mapping: Match AI-related patents to activities via patent text $\leftrightarrow$ activity text similarity. Independent of Recipe X geometry (which uses importance ratings, not text).
    \end{enumerate}
    \item Document which independence claim holds: independent of Recipe X (rating-cooccurrence)? Of Recipe Y (text embedding)? Of both?
\end{itemize}


\medskip
\noindent\textbf{PROFILE SELECTION PROTOCOL:}


\noindent\textit{Primary-profile approach (baseline):}
\begin{itemize}
    \item Designate $v_3$ as primary profile ex ante.
    \item $v_1$ and $v_2$ are robustness checks and hypothesis tests (``Does capability alone suffice? Does structure add?'').
    \item No profile selection based on outcome fit.
\end{itemize}


\noindent\textit{Alternative (if multiple $v_3$ candidates):}
\begin{itemize}
    \item Use holdout protocol: training window (2022--2023) for profile selection; test window (2024+) for headline results.
    \item Selection criterion: predictive $R^2$ for employment changes on training sample.
    \item Report test-window performance only as headline.
\end{itemize}


\medskip
\noindent\textbf{STABILITY TEST:}

Compare exposure rankings across shock types:
\begin{itemize}
    \item RBTC shock (routine-biased): $I^{\text{RBTC}}(a) = $ loading on routine activities
    \item AI shock (generative): $I^{\text{AI}}(a) = $ loading on language/cognitive activities
\end{itemize}

Compute rank correlation of $(E^{\text{RBTC}}_j, E^{\text{AI}}_j)$ across occupations. If correlation is low, different shocks have different exposure geometry; the framework may need shock-specific calibration rather than a single stable geometry.


\medskip
\noindent\textbf{DISCRIMINABILITY REQUIREMENT:}
\begin{itemize}
    \item Compute pairwise correlations among $E^{v_1}$, $E^{v_2}$, $E^{v_3}$ across occupations.
    \item If any pair has correlation $> 0.95$, they are empirically indistinguishable; note this and interpret comparisons cautiously.
\end{itemize}


\medskip
\noindent\textbf{OUTPUT:}
\begin{itemize}
    \item Discriminability table: correlations among candidate exposures across ($\mu$, domain) variants.
    \item Comparison table: fit statistics for $v_1$, $v_2$, $v_3$ vs discrete baselines (Eloundou et al.\ categorical; routine/manual indicators).
    \item Stability test: rank correlation of RBTC vs AI exposure rankings.
    \item Robustness table: stability of $v_1$/$v_2$/$v_3$ rankings across domain granularity.
\end{itemize}


\subsection{Model Evaluation Criteria}
\label{subsec:evaluation}


The framework is evaluated against simple baselines using explicit criteria.


\medskip
\noindent\textbf{PRIMARY CRITERIA (all must be addressed):}


\begin{enumerate}
    \item \textbf{Incremental explanatory power.} Activity-based exposure should improve fit relative to discrete baselines by an economically meaningful margin (e.g., $\Delta R^2 > 0.02$ or $\Delta$coefficient precision improvement). Evaluated out-of-sample (by time or held-out units). Phase~I establishes this criterion is met: normalized kernel overlap improves $R^2$ by 191\% over binary Jaccard.
    
    \item \textbf{Parameter stability.} Results should be qualitatively stable across domain choices (DWA vs GWA) and weighting schemes (uniform vs importance-weighted $\mu$).
    
    \item \textbf{Non-redundancy.} If correlation with discrete baselines $> 0.90$, value depends on:
    \begin{enumerate}
        \item[(a)] Interaction effects: Does activity-based exposure $\times$ wage-setting environment (union coverage, concentration, licensing) reveal heterogeneity that discrete baselines miss?
        \item[(b)] Out-of-sample gains: Even if in-sample correlation is high, does activity-based measure predict better out-of-sample?
    \end{enumerate}
    Phase~I shows correlation between binary Jaccard and semantic overlap is $r = 0.38$---substantial but not redundant.
    
    \item \textbf{Non-tautological validation.} Geometry diagnostics (Phase I) must use targets not used in construction. Phase I is rejection-based; passing Phase I does not validate the measure.
\end{enumerate}


\medskip
\noindent\textbf{PROFILE ADJUDICATION (Experiment B):}


\begin{enumerate}
    \item \textbf{Discriminability:} candidate exposures should have pairwise correlation $< 0.95$.
    
    \item \textbf{Incremental content:} at least one candidate should improve fit vs discrete baseline.
    
    \item \textbf{Consistency:} if $v_2 > v_1$, check whether the gap is larger in occupations with more structured work (per the ``structure matters'' hypothesis).
\end{enumerate}


\medskip
\noindent\textbf{REVISION PROTOCOL:}


If primary criteria fail:
\begin{itemize}
    \item Do not add auxiliary structure to rescue fit.
    \item Revise the shock profile construction (try different mapping) or domain granularity.
    \item If revision search is extensive, report null result: ``activity-based measure does not outperform discrete baselines under tested specifications.''
\end{itemize}


\section{Extensions, Limitations, and Historical Regimes}
\label{sec:extensions}


This paper targets a compact, implementable measurement pipeline. Extensions are recorded here for future work.


\subsection{Dynamic accumulation and reversals}
\label{subsec:dynamic}


The baseline construction is static. For settings where shock accumulation or reversal matters, define:
\[
A_t = \delta A_{t-1} + I_t + \varepsilon_t,
\]
where $\delta \in [0,1]$ governs persistence. This nests the static case ($\delta = 0$, single-period shock) and allows for:
\begin{itemize}
    \item \textbf{Accumulation ($\delta > 0$):} past shocks continue to affect current exposure.
    \item \textbf{Reversal ($\delta < 1$, negative $I_t$ shocks):} displacement can decline if technology retreats or comparative advantage shifts.
\end{itemize}


Estimation of the dynamic model requires panel structure at the occupation-time level and is deferred to future work.


\subsection{Augmentation channel}
\label{subsec:augmentation}


The framework carries $C_t$ (complementarity field) and $E^C$ (augmentation exposure) for symmetry. Estimation requires either:
\begin{itemize}
    \item Direct measurement of task-level productivity complementarities (not available in O*NET), or
    \item Auxiliary data on capital-labor complementarity by activity.
\end{itemize}


This channel is not estimated in the baseline empirics. A natural proxy is occupation-level measures of tool/technology adoption; validation would follow the same Phase I / Phase II structure as displacement exposure.


\subsection{Supervised dimension extraction}
\label{subsec:supervised}


An alternative to unsupervised text embeddings is supervised extraction of economically-relevant dimensions. Given theoretical priors about dimensions that matter (e.g., routine/non-routine, cognitive/manual), one can:
\begin{enumerate}
    \item Construct contrast sets of activities at each pole.
    \item Extract representations from a language model.
    \item Train linear probes to identify concept directions.
    \item Project all activities onto these dimensions.
\end{enumerate}

This approach directly extracts \emph{economic} structure rather than hoping unsupervised methods find it. The trade-off is that it requires theoretical priors about which dimensions matter, whereas unsupervised methods are agnostic. Phase~I results suggest that unsupervised text embeddings already capture economically relevant structure, but supervised approaches may improve further.


\subsection{Worker mobility as alternative validation target}
\label{subsec:mobility}


Wage comovement can be driven by industry concentration, geographic clustering, and institutional factors that operate orthogonal to activity structure. Worker mobility (occupation-to-occupation transitions) is more mechanically tied to ``can this worker do that job?''---the core question of activity overlap.

If kernel-weighted overlap predicts mobility as well as or better than it predicts wage comovement, the framework is robust across validation targets. This provides an additional test of whether the semantic structure captures economic substitutability.


\subsection{Alternative continuous representations}
\label{subsec:alternative-reps}


Phase~II representation comparison tested multiple approaches:

\begin{itemize}
    \item \textbf{O*NET structured dimensions.} Abilities (52 dimensions), Skills (35 dimensions), Knowledge (33 dimensions) provide occupation-level similarity measures that bypass activity-level text embeddings. These test a different hypothesis: ``Do occupations with similar ability/skill requirements comove?'' rather than ``Do occupations sharing semantically similar activities comove?''
    
    \item \textbf{Domain-specific embeddings.} JobBERT (trained on job postings) underperformed generic MPNet embeddings, suggesting domain-specific training does not improve prediction for this application.
    
    \item \textbf{Hybrid models.} Combining binary Jaccard with continuous measures (e.g., $H_2$: Jaccard + Abilities) achieves higher $R^2$ than either alone, suggesting the measures capture complementary information.
\end{itemize}

The activity-level text embedding approach (Recipe Y) remains the primary specification because it directly tests the task-manifold theory. O*NET structured dimensions, while predictive, test a different construct.


\subsection{Historical regime signatures}


To summarize historical episodes without forcing a single narrative, \cref{tab:regimes} records regime signatures as hypotheses about shock support, displacement vs augmentation, and primary mechanism.


\begin{table}[h!]
\centering
\begin{tabular}{|l|l|l|l|l|}
\hline
\textbf{Regime} & \textbf{Shock support ($I_t$)} & \textbf{$A$ vs.\ $C$} & \textbf{Mechanism} & \textbf{Wage wedge} \\
\hline
Electrification (1900--1940) & Manual/craft tasks & Mixed & TBD & Context-dependent \\
\hline
Postwar scaling (1945--1970) & Broad/diffuse & Augmentation & TBD & High pass-through \\
\hline
ICT/RBTC (1980--2007) & Routine-cognitive core & Displacement & Continuous & Declining pass-through \\
\hline
AI era (2020--) & Analytic/generative & TBD & TBD & TBD \\
\hline
\end{tabular}
\caption{Historical regimes as parameter signatures. ``Mechanism'' column records whether continuous or discrete activity structure is the primary channel; Phase~I suggests both are informative, with continuous explaining more variance and discrete providing higher precision.}
\label{tab:regimes}
\end{table}


\section{Conclusion}
\label{sec:conclusion}


This paper develops a geometric measurement framework for labor market exposure to technological shocks and tests whether continuous structure improves prediction beyond discrete activity classification. We model occupations as probability distributions over work activities, with shocks propagating through the activity space via a distance-based kernel. The framework provides a principled construction of occupation-activity measures from O*NET data and a testable formalization that distinguishes continuous from discrete representations of the task domain.


The central empirical finding is that both continuous and discrete representations are informative for predicting wage comovement, with neither unambiguously dominating. Normalized kernel-weighted overlap achieves $R^2 = 0.00485$---a 191\% improvement in explanatory power over binary Jaccard overlap ($R^2 = 0.00167$). However, binary Jaccard achieves higher statistical precision ($t = 8.00$ vs $t = 7.14$ for normalized kernel under clustered standard errors). The semantic signal is robust: it survives permutation testing ($p < 0.001$), remains significant after controlling for occupational concentration ($t = 5.29$ with entropy controls), and achieves the 100th percentile against random baseline distributions.


A critical methodological finding is that kernel bandwidth must be calibrated to local distance structure. Setting $\sigma$ to the global distance median causes kernel collapse, making continuous measures equivalent to noise---an artifact that initially led to an incorrect conclusion that ``discrete dominates continuous.'' Calibrating $\sigma$ to the median nearest-neighbor distance ($\sigma = 0.223$ for our embeddings) restores discrimination and reveals the continuous structure.


Normalization for concentration effects improves prediction. The normalized overlap measure, which controls for specialist-vs-generalist distributional properties, achieves 56.6\% higher $R^2$ than unnormalized overlap. This indicates that concentration was noise obscuring the semantic signal: wage comovement is driven by the semantic content of shared activities, not by whether occupations are similarly specialized or diffuse.


Several questions remain for Phase~II. First, we have not tested whether activity-based exposure is the \emph{dominant} channel versus industry composition, credential similarity, or wage proximity. The dominance test proposed in Phase~II will assess this. Second, we have not validated structural stability across shock types; the RBTC-vs-AI stability test is needed. Third, wage comovement may not be the ideal validation target---worker mobility may provide a more direct test of economic substitutability. Fourth, the framework has validated the propagation geometry but not the shock profiles; constructing and validating AI-specific profiles ($v_1$, $v_2$, $v_3$) is the critical next step.


The framework's contribution is a disciplined measurement pipeline that reveals what kind of structure actually matters for economic outcomes. The finding that both continuous and discrete representations are informative---with continuous explaining more variance and discrete achieving higher precision---suggests that the choice depends on research priorities. Continuous kernel-weighted measures capture economically meaningful variation that binary classification misses, while discrete measures provide more robust inference. The methodological lessons---calibrate bandwidth to local structure, normalize for concentration, validate against random baselines---provide guidance for future implementations.


\bibliographystyle{apalike}
\bibliography{references}


\end{document}